\documentclass[../DoAn.tex]{subfiles}
\begin{document}

Dựa trên cấu trúc thiết kế kịch bản và mỹ thuật được đặc tả trong Chương trước, chương này tập trung phân tích thiết kế chi tiết ở mức mã nguồn và kiến trúc phần mềm của trò chơi "Dream Witch". Hệ thống được lập trình theo mô hình hướng thành phần (Component-Based Architecture) của Unity Engine và áp dụng các mẫu thiết kế (Design Patterns) phổ biến nhằm đảm bảo tính mở rộng, khả năng bảo trì và hiệu năng tối ưu trên các thiết bị.

\section{Thiết kế kiến trúc hệ thống và phân rã thành phần}
\label{section:4.1}

Hệ thống mã nguồn của trò chơi được tổ chức dựa trên việc kết hợp giữa kiến trúc hướng thành phần (Component-Based Architecture) đặc thù của Unity Engine và các nguyên lý thiết kế phần mềm hướng kiến trúc phân tầng. Sơ đồ kiến trúc tổng thể các lớp logic và dòng dữ liệu tương tác giữa các phân hệ được mô tả trong Hình~\ref{fig:game_architecture}.

\begin{figure}[H]
    \centering
    \includegraphics[width=0.65\linewidth,height=0.22\textheight,keepaspectratio]{Hinhve/game_architecture.png}
    \caption{Kiến trúc tổng thể các phân hệ và lớp logic trong hệ thống}
    \label{fig:game_architecture}
\end{figure}

\subsection{Lựa chọn kiến trúc phần mềm}
\label{subsection:4.1.1}
Để giải quyết bài toán quản lý trạng thái trò chơi phức tạp của thể loại Metroidvania, đề tài lựa chọn áp dụng mô hình kiến trúc MVC (Model-View-Controller) kết hợp với mô hình hướng thành phần của Unity:
\begin{itemize}
    \item \textbf{Model (Mô hình dữ liệu):} Chịu trách nhiệm quản lý, lưu trữ trạng thái chạy và dữ liệu tĩnh của trò chơi. Trong dự án, thành phần này được hiện thực hóa thông qua các đối tượng \texttt{ScriptableObject} bền vững như:
    \begin{itemize}
        \item \texttt{PlayerFormDataSO:} Lưu cấu hình tĩnh cho 6 hình thái linh hồn (HP tối đa, Animator Controller, Prefab vật lý).
        \item \texttt{ItemDatabaseSO} và \texttt{SpellDatabaseSO:} Định nghĩa danh sách các vật phẩm và phép thuật.
        \item \texttt{InventoryStateSO}, \texttt{ShopStateSO}, \texttt{CurrencyWalletSO:} Lưu trạng thái runtime động của túi đồ, cửa hàng và số dư tiền của nhân vật.
    \end{itemize}
    \item \textbf{View (Giao diện hiển thị):} Chịu trách nhiệm hiển thị thông tin trực quan lên màn hình và ghi nhận thao tác của người chơi để chuyển giao cho Controller xử lý. Thành phần View bao gồm các đối tượng Canvas UI tĩnh và động như \texttt{PlayerHudUI} (HUD hiển thị HP/Stamina/Mana), các panel quản lý cửa hàng (\texttt{ShopUI}), túi đồ (\texttt{InventoryUI}), nâng cấp (\texttt{FacilityUI}), và chữ hiển thị sát thương bay (\texttt{DamageTextView}).
    \item \textbf{Controller (Bộ điều phối):} Đóng vai trò trung gian xử lý các logic nghiệp vụ, lắng nghe tương tác từ View để thay đổi Model, và ngược lại. Các controller tiêu biểu là:
    \begin{itemize}
        \item \texttt{PlayerController:} Điều phối toàn bộ trạng thái di chuyển và chiến đấu của người chơi.
        \item \texttt{UIManager} và \texttt{UIStateManager:} Quản lý vòng đời và thứ tự ưu tiên đóng mở các Canvas UI.
        \item \texttt{GameSaveManager:} Đọc/ghi trạng thái từ các lớp Model xuống tệp tin lưu trữ JSON.
    \end{itemize}
\end{itemize}

\subsection{Thiết kế tổng quan (Biểu đồ gói)}
\label{subsection:4.1.2}
Kiến trúc phần mềm của trò chơi được phân chia thành các gói (packages) và được tổ chức theo mô hình phân tầng chặt chẽ từ trên xuống dưới, đảm bảo tính độc lập và khả năng bảo trì cao:
\begin{itemize}
    \item \textbf{Tầng Giao diện (View Layer - Gói UI):} Chứa các lớp hiển thị giao diện người dùng. Tầng này chỉ phụ thuộc vào tầng điều phối bên dưới để lấy dữ liệu hiển thị, không chứa logic xử lý vật lý hay dữ liệu lưu trữ trực tiếp.
    \item \textbf{Tầng Điều phối và Logic (Controller/Logic Layer - Gói Player, Enemy, System):} Chứa logic nghiệp vụ cốt lõi của trò chơi bao gồm bộ điều khiển nhân vật chính, AI quái vật tuần tra/Boss, và các hệ thống dịch vụ nền như Save/Load, Culling, Pause, Shop.
    \item \textbf{Tầng Hạ tầng và Dữ liệu (Model/Utility Layer - Gói Core, Interfaces):} Định nghĩa các hằng số cấu hình hệ thống, mẫu Singleton dùng chung (\texttt{SingleInstance}), và các giao diện trừu tượng (\texttt{IDamageable}, \texttt{IInteractable}, \texttt{ICullable}). Tầng này hoàn toàn độc lập và không phụ thuộc ngược lên các tầng trên.
\end{itemize}

Để đảm bảo quy tắc thiết kế không tạo phụ thuộc vòng (circular dependency) giữa gói \texttt{Player} và \texttt{Enemy} trong tầng Logic (khi nhân vật cần phát hiện va chạm gây sát thương lên quái vật và quái vật cần dò tìm vị trí nhân vật), cả hai gói đều giao tiếp thông qua các giao diện trừu tượng được định nghĩa sẵn ở gói \texttt{Interfaces} như \texttt{IDamageable}. Lớp \texttt{PlayerController} và \texttt{MvEnemyBase} đều thực thi giao diện này, giúp giảm thiểu độ kết hợp cứng (tight coupling) giữa hai phân hệ.

Sơ đồ phân rã các phân hệ chính trong mã nguồn trò chơi theo cấu trúc thư mục được đặc tả trong Hình~\ref{fig:SoDoPhanRaPhanHe}.

\begin{figure}[H]
    \centering
    \includegraphics[width=0.65\linewidth,height=0.22\textheight,keepaspectratio]{Hinhve/system_decomposition.png}
    \caption{Sơ đồ phân rã các phân hệ chính trong mã nguồn trò chơi}
    \label{fig:SoDoPhanRaPhanHe}
\end{figure}

\subsection{Thiết kế chi tiết gói (Biểu đồ lớp)}
\label{subsection:4.1.3}
Trong thiết kế chi tiết gói, mối quan hệ tương tác giữa các lớp bên trong gói \texttt{Player} và gói \texttt{Enemy} tuân thủ chặt chẽ các quan hệ UML tiêu chuẩn:
\begin{itemize}
    \item \textbf{Quan hệ Hợp thành (Composition):} Lớp điều phối \texttt{PlayerController} sở hữu máy trạng thái \texttt{PlayerStateMachine} nhằm quản lý vòng đời và các hành vi chuyển động độc lập của nhân vật. Khi thực thể Player bị hủy bỏ khỏi bộ nhớ, máy trạng thái cũng sẽ bị giải phóng hoàn toàn theo.
    \item \textbf{Quan hệ Kết hợp (Association):} Lớp \texttt{PlayerController} duy trì các liên kết đến các thành phần con độc lập khác như \texttt{PlayerMovement} (xử lý vật lý di chuyển), \texttt{PlayerCombat} (xử lý tấn công), \texttt{PlayerStats} (chỉ số sinh tồn), và \texttt{PlayerFormManager} (quản lý biến hình) để ủy quyền (delegate) thực hiện các hành vi tương ứng.
    \item \textbf{Quan hệ Kết tập (Aggregation):} Lớp \texttt{PlayerFormManager} chứa một danh sách tập hợp các cấu hình linh hồn \texttt{PlayerFormDataSO}. Các ScriptableObjects cấu hình này tồn tại độc lập dưới dạng tài nguyên lưu trữ trong Asset Database của Unity và không bị hủy khi lớp quản lý bị hủy.
    \item \textbf{Quan hệ Phụ thuộc (Dependency):} Lớp \texttt{PlayerMovement} phụ thuộc vào lớp tham chiếu trung gian \texttt{PlayerFormBodyRef} để lấy các thông số collider va chạm và vị trí các điểm kiểm tra mặt đất, vách tường một cách động mỗi khi người chơi thay đổi hình thái linh hồn.
    \item \textbf{Quan hệ Kế thừa (Inheritance):} Các lớp Enemy cụ thể (như \texttt{MvEm0010} - BoneDog) kế thừa từ lớp cơ sở chung là \texttt{MvEnemyBase} để tái sử dụng toàn bộ cấu trúc máy trạng thái quái vật, hệ thống tuần tra vật lý và logic nhận sát thương.
    \item \textbf{Quan hệ Thực thi (Implementation):} Lớp \texttt{PlayerController} và lớp cơ sở \texttt{MvEnemyBase} đều thực thi giao diện \texttt{IDamageable} của gói \texttt{Interfaces} nhằm đáp ứng cơ chế nhận sát thương từ các nguồn đòn đánh khác nhau.
\end{itemize}


\section{Phân hệ Người chơi và Cơ chế Biến hình (Form Switch)}
\label{section:4.2}

\subsection{Thiết kế dữ liệu và Prefab cho các hình thái linh hồn (Player Forms)}
Cơ chế biến hình cho phép người chơi thay đổi linh hoạt hình dạng của nhân vật chính sang các dạng quái vật khác khi thu thập linh hồn. Hệ thống này sử dụng cấu trúc dữ liệu tĩnh hướng cấu hình (Data-Driven Design) thông qua các lớp:
\begin{itemize}
    \item \textbf{\texttt{PlayerFormId} (Enum):} Định danh các hình thái: \texttt{Human} (Dạng người mặc định), \texttt{Em0010} (Chó xương - BoneDog), \texttt{Em0020} (Mèo ma - GhostCat), \texttt{Em0060} (Hiệp sĩ khiên - Shield), \texttt{Em0070} (Phù thủy bay - Pile), và \texttt{Em0100} (Đấu sĩ - Minotaur).
    \item \textbf{\texttt{PlayerFormDataSO} (ScriptableObject):} Lưu trữ các dữ liệu thiết kế cố định như tên hiển thị, icon HUD, lượng HP tối đa, Animator Controller và Prefab vật lý tương ứng của từng hình thái.
    \item \textbf{\texttt{PlayerFormConfig}:} Quản lý việc mở khóa các form và lưu trữ chỉ số sinh lực thực tế khi đang chạy (Runtime Health) của từng linh hồn độc lập.
\end{itemize}

\subsection{Quy trình đồng bộ hóa chuyển đổi hình thái vật lý và hoạt ảnh}
Khi người chơi kích hoạt chuyển đổi linh hồn gần một xác chết kẻ địch, lớp điều phối trung tâm \texttt{PlayerController} kết hợp với \texttt{PlayerFormManager} thực hiện chuỗi hành vi đồng bộ hóa thời gian thực:
\begin{enumerate}
    \item \textbf{Đồng bộ hóa Chỉ số Sinh lực:} Lưu chỉ số HP hiện tại của form cũ và thiết lập chỉ số HP tương ứng của form mới vào lớp \texttt{PlayerStats}.
    \item \textbf{Tải cấu trúc vật lý động:} Lớp \texttt{PlayerFormManager} hủy phiên bản Prefab cũ và khởi tạo (Instantiate) Prefab của form mới làm con của Player Root. Prefab này được đính kèm component \texttt{PlayerFormBodyRef} định nghĩa rõ ràng các collider va chạm môi trường và các trigger hitbox tấn công cận chiến đặc thù.
    \item \textbf{Đồng bộ hóa Collider trong Physics Engine:} Lớp \texttt{PlayerMovement} tiếp nhận thông tin từ \texttt{PlayerFormBodyRef} mới để gán lại Collider va chạm thực tế và các điểm kiểm tra mặt đất (\texttt{GroundCheck}), vách tường (\texttt{WallCheck}) vào Physics2D của Unity, tránh lỗi kẹt địa hình do chênh lệch kích thước giữa các hình thái.
    \item \textbf{Rebind hoạt ảnh:} Thay đổi runtime controller của thành phần \texttt{Animator} trên Player Root bằng Animator Controller của form mới để kích hoạt tức thì các clip hoạt ảnh tương ứng. Hệ thống quản lý hoạt ảnh của nhân vật chính (\texttt{PlayerAnimationController}) và các quái vật thông thường sẽ thực hiện gọi trực tiếp các hoạt ảnh này thông qua chuỗi ký tự (string-based loading) thay vì sử dụng các tham số và quan hệ chuyển cảnh phức tạp trong Animator Controller.
    \item \textbf{Cập nhật chỉ số chiến đấu:} Thay đổi cấu hình sát thương và chuỗi combo cận chiến trong \texttt{PlayerCombat} theo thông số của form mới.
\end{enumerate}

Sơ đồ trình tự biểu diễn quá trình đồng bộ hóa khi chuyển đổi hình thái linh hồn được mô tả trong Hình~\ref{fig:form_switch_sequence}.

\begin{figure}[H]
    \centering
    \includegraphics[width=0.9\linewidth,height=0.25\textheight,keepaspectratio]{Hinhve/FormSwitchSequence.png}
    \caption{Sơ đồ trình tự quá trình đồng bộ hóa chuyển đổi hình thái (Form Switch)}
    \label{fig:form_switch_sequence}
\end{figure}
% NOTE: Sinh viên có thể chèn hình chụp hiệu ứng hình thái linh hồn (Form Switch Intro VFX) thực tế tại đây
 \begin{figure}[H]
    \centering
    \includegraphics[width=0.75\linewidth]{Hinhve/BienHinh.png}
    \caption{Interact Prompt hiện chuyển đổi linh hồn (Form Switch)}
    \label{fig:form_switch_vfx}
 \end{figure}


Quy trình thực thi của phương thức \texttt{ActivateForm} trong lớp \texttt{PlayerFormManager} nhằm quản lý vòng đời và thiết lập các thông số vật lý cho hình thái linh hồn mới được thực hiện theo các bước sau:
\begin{itemize}
    \item \textbf{Bước 1: Kiểm tra dữ liệu đầu vào:} Kiểm tra cấu hình hình thái mới (\texttt{PlayerFormDataSO}). Nếu dữ liệu rỗng (\texttt{null}), phương thức sẽ lập tức kết thúc hành vi để tránh lỗi runtime.
    \item \textbf{Bước 2: Hủy/Vô hiệu hóa hình thái hiện tại:} Nếu có một hình thái linh hồn khác đang hoạt động, hệ thống sẽ tiến hành hủy hoặc đưa về trạng thái vô hiệu hóa. Đồng thời, đối tượng gốc của dạng người mặc định (\texttt{humanFormRoot}) cũng sẽ được ẩn đi.
    \item \textbf{Bước 3: Khởi tạo Prefab vật lý mới:} Khởi tạo (\texttt{Instantiate}) Prefab tương ứng của hình thái linh hồn mới làm đối tượng con trực tiếp của Player Root, thiết lập các thuộc tính tọa độ, góc quay và tỷ lệ kích thước về giá trị mặc định (\texttt{zero}/\texttt{identity}/\texttt{one}).
    \item \textbf{Bước 4: Trích xuất tham chiếu cấu trúc vật lý:} Truy vấn thành phần \texttt{PlayerFormBodyRef} nằm trong đối tượng mới để lấy tham chiếu đến bộ Collider va chạm (\texttt{BodyCollider}), làm cơ sở đồng bộ hóa với hệ thống Physics2D.
    \item \textbf{Bước 5: Thiết lập trạng thái phòng thủ:} Vô hiệu hóa Collider bảo vệ (\texttt{GuardCollider}) để chuẩn bị cho chu kỳ cập nhật trạng thái mới của nhân vật.
\end{itemize}

\subsection{Quản lý trạng thái thông qua máy trạng thái động (Form State Factory)}
Để tránh việc lớp \texttt{PlayerController} bị phình to do phải quản lý hàng chục trạng thái của nhiều form khác nhau, hệ thống sử dụng mẫu thiết kế Factory thông qua lớp \texttt{FormStateFactory}. 

Khi người chơi ở dạng người mặc định, \texttt{PlayerStateMachine} sử dụng các trạng thái tĩnh được khởi tạo từ đầu (như \texttt{IdleState}, \texttt{MoveState}, \texttt{JumpState}, \texttt{DashState}, \texttt{WallClimbState}). Khi chuyển sang một dạng Enemy biến hình, \texttt{FormStateFactory} sẽ tạo ra ``on-demand'' (theo nhu cầu) một bộ trạng thái chuyên biệt cho form đó (ví dụ: \texttt{Em0010IdleState}, \texttt{Em0010AttackState}) kế thừa từ lớp cơ sở \texttt{PlayerState}.

\section{Phân hệ Kẻ địch và Trí tuệ nhân tạo (AI Brain)}
\label{section:4.3}

\subsection{Phân cấp kế thừa các thực thể (Enemy Hierarchy)}
Hệ thống quái vật trong game được thiết kế hướng đối tượng chặt chẽ nhằm tối ưu hóa việc quản lý mã nguồn di chuyển, phát hiện mục tiêu và va chạm vật lý:
\begin{itemize}
    \item \textbf{Lớp cơ sở trừu tượng \texttt{MvEnemyBase}:} Kế thừa trực tiếp từ \texttt{MonoBehaviour} và thực thi giao diện \texttt{IDamageable}. Lớp này chịu trách nhiệm quản lý máy trạng thái \texttt{EnemyStateMachine}, quản lý các chỉ số cơ bản (Máu, Tốc độ, Tầm đánh), xử lý các hiệu ứng vật lý (lực đẩy lùi - Knockback, đệm va chạm), và tích hợp hệ thống tự động tìm kiếm mục tiêu (\texttt{UpdateSearchTarget}) trong phạm vi quạt góc quét (\texttt{searchFanAngle}) và tầm quét (\texttt{searchRadius}). Đồng thời, lớp này khai báo cấu trúc enum \texttt{AsCommon} đại diện cho các trạng thái dùng chung của quái vật:
    \begin{itemize}
        \item Các trạng thái di chuyển và tuần tra: \texttt{Idle} (0), \texttt{Run} (1), \texttt{Turn} (2), và \texttt{AtkAfter} (3).
        \item Các trạng thái nhận phản hồi sát thương: \texttt{Hit} (4), \texttt{HitKnockBack} (5), \texttt{HitKnockUp} (6), \texttt{HitBounce} (7), và \texttt{HitNeedle} (8).
        \item Các trạng thái bất lợi: \texttt{CondFrozen} (9), \texttt{CondStone} (10), và \texttt{CondStun} (11).
        \item Các trạng thái vòng đời và kết thúc: \texttt{DeathPotDive} (12), \texttt{DeathPre} (13), \texttt{Cage} (14), \texttt{Death} (15), \texttt{Inactive} (16), và \texttt{Max} (17).
    \end{itemize}
    \item \textbf{Các lớp quái vật và Boss cụ thể (\texttt{MvEmXXXX}):} Kế thừa từ \texttt{MvEnemyBase} (ví dụ: \texttt{MvEm0010} - BoneDog, \texttt{MvEm9020} - Golem Boss, \texttt{MvEm9030} - Marionette Boss). Các lớp này override các phương thức đặc thù như thiết lập kiểu quái vật (\texttt{EmType}), ánh xạ tên trạng thái (\texttt{getAsEnumName\_EmUnique}), đăng ký bảng trạng thái (\texttt{makeAsTable}), và tùy biến thuật toán quét tìm người chơi (\texttt{moveSearchPl}). Các lớp này cũng định nghĩa tập hợp các trạng thái riêng biệt thông qua enum \texttt{As} bắt đầu từ chỉ mục 17 (ví dụ BoneDog có thêm \texttt{AtkCharge} = 17, \texttt{Chase} = 19, \texttt{Jump} = 20).
\end{itemize}

\begin{figure}[H]
    \centering
    \includegraphics[width=0.9\linewidth,height=0.25\textheight,keepaspectratio]{Hinhve/enemy_class_hierarchy.png}
    \caption{Biểu đồ phân cấp kế thừa các thực thể Kẻ địch (Enemy Class Hierarchy)}
    \label{fig:EnemyClassHierarchy}
\end{figure}


\subsection{Cơ chế AI cho Boss (Brain, Desire và Combo)}
Các trận chiến với Boss lớn (như Big Golem \texttt{Em9020} hay Marionette \texttt{Em9030}) được điều phối bởi lớp bộ não AI chuyên biệt \texttt{MvEmBrain} chạy song song với máy trạng thái:
\begin{itemize}
    \item \textbf{Trọng số mong muốn (Desire Weights):} AI liên tục đánh giá khoảng cách đến người chơi và phần trăm máu còn lại của Boss để cập nhật bảng trọng số mong muốn của các chiêu thức (ví dụ: khi người chơi ở xa, trọng số chiêu phóng đạn tăng; khi Boss dưới 50\% máu, trọng số các đòn tấn công nhanh, diện rộng tăng).
    \item \textbf{Chuỗi chiêu thức (Combo Patterns):} Thay vì ra đòn ngẫu nhiên, Boss chọn một chuỗi hành động kết hợp được lập trình sẵn (\texttt{ComboBuff} lưu mảng các trạng thái hành vi). Khi combo bắt đầu, Boss sẽ thực thi tuần tự tất cả các đòn đánh trong combo đó (ví dụ: Tiếp cận $\rightarrow$ Đập đất $\rightarrow$ Phóng tia năng lượng) trước khi quay lại trạng thái đánh giá mong muốn tiếp theo.
\end{itemize}

Sơ đồ luồng xử lý quyết định hành động của Boss AI được mô phỏng trong Hình~\ref{fig:boss_ai_flow}.

\begin{figure}[H]
    \centering
    \includegraphics[width=0.9\linewidth,height=0.25\textheight,keepaspectratio]{Hinhve/BossAiFlow.png}
    \caption{Sơ đồ luồng ra quyết định hành động của Boss AI dựa trên Desire và Combo}
    \label{fig:boss_ai_flow}
\end{figure}

\section{Các hệ thống lõi phục vụ vận hành}
\label{section:4.4}

\subsection{Hệ thống Lưu trữ và Phục hồi dữ liệu (Save/Load)}
\label{section:4.4.1}
Hệ thống sử dụng lớp \texttt{GameSaveManager} để serialize trạng thái trò chơi thành chuỗi dữ liệu JSON và lưu trữ lâu dài:
\begin{itemize}
    \item \textbf{Quản lý dữ liệu người chơi:} Ghi lại lượng tiền vàng (\texttt{gold}), danh sách vật phẩm trong hòm đồ (\texttt{inventory}), bình máu và công cụ đang trang bị.
    \item \textbf{Quản lý trạng thái tiến trình:} Lưu trữ các nâng cấp vĩnh viễn đã mua tại làng, các kỹ năng phép thuật đã học, danh sách các cổng dịch chuyển (\texttt{Portals}) đã kích hoạt, các cửa lớn đã mở khóa và danh sách các Boss đã bị tiêu diệt.
    \item \textbf{Độc lập cảnh chơi:} \texttt{GameSaveManager} được thiết kế theo mẫu Singleton và nằm trên Object có thuộc tính \texttt{DontDestroyOnLoad}, đảm bảo tiến trình lưu/tải không bị gián đoạn khi chuyển cảnh.
\end{itemize}

Hệ thống quản lý lưu trữ được phân tách thành hai dạng độc lập: lưu trữ cục bộ các thiết lập người dùng bằng cơ chế \texttt{PlayerPrefs} của Unity, và lưu trữ tiến trình chơi game bằng tệp tin định dạng JSON bền vững.

Bảng~\ref{table:player_prefs_storage} đặc tả các khóa thiết lập cấu hình hệ thống được lưu trữ thông qua \texttt{PlayerPrefs}.

\begin{table}[H]
\centering
\begin{tabular}{|p{4.5cm}|p{2.0cm}|p{2.0cm}|p{5.5cm}|}
\hline
\textbf{Khóa lưu trữ (Key)} & \textbf{Kiểu dữ liệu} & \textbf{Giá trị mặc định} & \textbf{Ý nghĩa chức năng} \\ \hline
\texttt{Pref\_MasterVolume} & float & 1.0f & Âm lượng tổng của trò chơi (Master Volume). \\ \hline
\texttt{Pref\_BgmVolume} & float & 0.8f & Âm lượng nhạc nền (BGM Volume). \\ \hline
\texttt{Pref\_SfxVolume} & float & 0.8f & Âm lượng hiệu ứng âm thanh (SFX Volume). \\ \hline
\end{tabular}
\caption{Bảng đặc tả các khóa lưu trữ cấu hình trong PlayerPrefs}
\label{table:player_prefs_storage}
\end{table}

Đối với tiến trình chơi game chi tiết, tệp tin cấu trúc JSON được ghi dưới dạng các slot lưu trữ độc lập tại đường dẫn \texttt{Application.persistentDataPath}. Các trường dữ liệu chính được lưu trong lớp \texttt{GameSaveData} được đặc tả chi tiết trong Bảng~\ref{table:json_save_data}.

\begin{table}[H]
\centering
\begin{tabular}{|p{4.0cm}|p{3.5cm}|p{6.5cm}|}
\hline
\textbf{Trường dữ liệu (JSON Field)} & \textbf{Kiểu dữ liệu C\#} & \textbf{Ý nghĩa lưu trữ thực tế} \\ \hline
\texttt{gold} & int & Số lượng tiền vàng tích lũy hiện có. \\ \hline
\texttt{inventory} & InventorySaveData & Danh sách các vật phẩm thu thập được trong túi đồ. \\ \hline
\texttt{toolEquip} & EquipmentSaveData & Dữ liệu về công cụ (Tool) đang trang bị. \\ \hline
\texttt{healingPotionEquip} & EquipmentSaveData & Dữ liệu về bình thuốc hồi phục đang trang bị. \\ \hline
\texttt{shop} & ShopSaveData & Trạng thái đã mua/chưa mua các vật phẩm tại cửa hàng. \\ \hline
\texttt{skillProgress} & SkillProgressSaveData & Trạng thái mở khóa các phép thuật đã học. \\ \hline
\texttt{facilityProgress} & FacilityProgressSaveData & Tiến trình nâng cấp vĩnh viễn các thuộc tính tại Village. \\ \hline
\texttt{skillTreeProgress} & SkillTreeProgressSaveData & Tiến trình mở khóa nhánh kỹ năng chiến đấu (Skill Tree). \\ \hline
\texttt{portals} & PortalSaveData & Danh sách các cổng dịch chuyển (Portals) đã kích hoạt. \\ \hline
\texttt{doors} & DoorSaveData & Trạng thái mở khóa các cánh cửa lớn trên thế giới. \\ \hline
\texttt{worldPickups} & WorldPickupSaveData & Trạng thái nhặt các vật phẩm thu thập trong màn chơi. \\ \hline
\texttt{bossDefeats} & BossDefeatSaveData & Trạng thái đã đánh bại Big Golem và Marionette. \\ \hline
\texttt{equippedSpellId} & string & ID của phép thuật đang được chọn. \\ \hline
\texttt{slotIndex} & int & Chỉ mục ô nhớ lưu trữ game (Slot Index). \\ \hline
\texttt{playTimeSeconds} & float & Tổng thời gian chơi game đã ghi nhận (tính bằng giây). \\ \hline
\end{tabular}
\caption{Bảng đặc tả cấu trúc lưu trữ tiến trình game trong GameSaveData}
\label{table:json_save_data}
\end{table}

Ngoài ra, để tối ưu hóa thiết kế hướng dữ liệu (Data-Driven Design), game sử dụng các ScriptableObjects được chia làm hai nhóm chính: dữ liệu cấu hình tĩnh và dữ liệu lưu trạng thái động. Vai trò và phân lớp cụ thể được liệt kê trong Bảng~\ref{table:scriptable_objects}.

\begin{table}[H]
\centering
\begin{tabular}{|p{4.5cm}|p{3.5cm}|p{6.0cm}|}
\hline
\textbf{Tên ScriptableObject} & \textbf{Lớp C\# kế thừa} & \textbf{Vai trò và Chức năng lưu trữ} \\ \hline
\texttt{ItemDatabaseSO} & ScriptableObject & Danh sách định nghĩa tĩnh toàn bộ vật phẩm trong game. \\ \hline
\texttt{SpellDatabaseSO} & ScriptableObject & Danh sách định nghĩa các phép thuật có thể học. \\ \hline
\texttt{PlayerFormDataSO} & ScriptableObject & Chứa cấu hình prefab, icon, máu mặc định của 6 form linh hồn. \\ \hline
\texttt{InventoryStateSO} & ScriptableObject & Runtime State: Quản lý danh sách vật phẩm chạy trong túi đồ. \\ \hline
\texttt{CurrencyWalletSO} & ScriptableObject & Runtime State: Quản lý số dư tiền vàng hiện thời của Player. \\ \hline
\texttt{ShopStateSO} & ScriptableObject & Runtime State: Theo dõi tiến trình mua sắm vật phẩm của shop. \\ \hline
\end{tabular}
\caption{Bảng phân loại các ScriptableObjects trong kiến trúc game}
\label{table:scriptable_objects}
\end{table}

Quy trình thực thi của phương thức \texttt{SaveGameToPath} trong lớp \texttt{GameSaveManager} thực hiện việc thu thập dữ liệu trạng thái trò chơi, tính toán thời gian chơi, tuần tự hóa sang định dạng JSON và ghi xuống tệp lưu trữ được tổ chức theo các bước logic sau:
\begin{itemize}
    \item \textbf{Bước 1: Nạp dữ liệu cũ (nếu có):} Tiến hành đọc tệp tin lưu trữ cũ tại đường dẫn chỉ định (\texttt{LoadDataFromPath}) để trích xuất thời điểm khởi tạo tệp game (\texttt{createdUtcTicks}) và thời gian chơi tích lũy trước đó.
    \item \textbf{Bước 2: Thu thập dữ liệu hiện thời:} Gọi phương thức \texttt{CaptureSaveData()} để tổng hợp toàn bộ trạng thái chạy hiện tại của nhân vật, ví đồ, tiến trình nâng cấp và lịch sử đánh bại Boss vào cấu trúc \texttt{GameSaveData}, đồng thời gán chỉ mục ô nhớ tương ứng (\texttt{slotIndex}).
    \item \textbf{Bước 3: Tính toán thời gian chơi và mốc thời gian (Timestamp):}
    \begin{itemize}
        \item Thiết lập mốc thời gian cập nhật tệp tin (\texttt{updatedUtcTicks}) theo thời gian hệ thống hiện hành.
        \item Nếu là tệp mới, gán mốc thời gian tạo bằng thời gian hiện tại.
        \item Tính toán tổng thời gian chơi lũy kế bằng cách cộng thêm thời gian chạy thực tế của phiên chơi hiện tại (\texttt{Time.unscaledTime}).
    \end{itemize}
    \item \textbf{Bước 4: Tuần tự hóa dữ liệu sang JSON:} Sử dụng thư viện \texttt{JsonUtility} của Unity để chuyển đổi đối tượng lưu trữ \texttt{GameSaveData} thành một chuỗi định dạng văn bản JSON có định dạng rõ ràng (pretty print) nhằm dễ đọc và gỡ lỗi.
    \item \textbf{Bước 5: Tạo thư mục và ghi tệp tin:} Kiểm tra sự tồn tại của thư mục đích trên ổ đĩa. Nếu chưa tồn tại, tiến hành tạo thư mục mới trước khi ghi toàn bộ nội dung chuỗi JSON xuống tệp tin bền vững bằng phương thức ghi luồng đồng bộ.
\end{itemize}
% NOTE: Sinh viên có thể chèn hình chụp thực tế đạn phép hoặc combat (Object Pooling) tại đây
 \begin{figure}[H]
    \centering
    \includegraphics[width=0.85\linewidth]{Hinhve/SpellGameplay.png}
    \caption{Hình ảnh đạn phép bay và combo vật thể được quản lý bởi Object Pooling}
    \label{fig:spell_pooling_screenshot}
 \end{figure}


\subsection{Hệ thống Culling tối ưu hóa hiệu năng (Culling System)}
Để duy trì tốc độ khung hình cao ổn định trên thiết bị phần cứng hạn chế, trò chơi tích hợp \texttt{CullingManager} quản lý việc bật tắt động các tài nguyên:
\begin{itemize}
    \item Bản đồ được chia thành các khu vực phòng độc lập (\texttt{RoomController}).
    \item Khi người chơi di chuyển ra khỏi phạm vi hoạt động của một phòng chơi, \texttt{CullingManager} sẽ tự động vô hiệu hóa (\texttt{SetActive(false)}) toàn bộ Sprite Renderers, Collision Colliders và tắt các script tính toán AI của quái vật thuộc phòng đó.
    \item Các đối tượng sẽ được kích hoạt lại ngay lập tức khi người chơi bước vào vùng biên giới tiếp giáp phòng, giúp giảm thiểu đáng kể số lượng lệnh vẽ (Draw Calls) và các phép tính toán vật lý không cần thiết.
\end{itemize}
% NOTE: Sinh viên có thể chèn hình chụp minh họa hệ thống Culling (ẩn các room bên ngoài) trong Unity Editor tại đây
 \begin{figure}[H]
    \centering
    \includegraphics[width=0.9\linewidth,height=0.25\textheight,keepaspectratio]{Hinhve/CullingVisualization.png}
    \caption{Hình vẽ trực quan các phòng chơi bị tắt hiển thị ngoài tầm nhìn để tối ưu hóa}
    \label{fig:culling_vis}
 \end{figure}


\subsection{Hệ thống Nâng cấp (Facility) và Cửa hàng (Shop)}
Hệ thống kinh tế và thăng tiến sức mạnh của nhân vật được vận hành bởi hai lớp:
\begin{itemize}
    \item \textbf{\texttt{FacilityManager}:} Đọc ghi tiến trình nâng cấp chỉ số vĩnh viễn (HP, Mana, Sát thương cận chiến, Sức chứa kho) dựa trên số lượng các vật phẩm đặc biệt mà người chơi thu thập được từ thế giới.
    \item \textbf{\texttt{ShopStateSO}:} Sử dụng ScriptableObject để lưu trữ danh sách các kỹ năng phép thuật và công cụ bổ trợ có bán tại NPC Shopkeeper, tự động đồng bộ hóa trạng thái đã mua/chưa mua giữa các màn chơi.
\end{itemize}

\section{Kiến trúc giao diện người dùng và Canvas (UI Canvas Architecture)}
\label{section:4.5}

Hệ thống giao diện người dùng (UI) trong "Dream Witch" được thiết kế dựa trên cấu trúc phân lớp Canvas của Unity Engine để phân tách hiển thị và tối ưu hóa hiệu năng render:
\begin{itemize}
    \item \textbf{HUD Canvas (Screen Space - Overlay):} Quản lý các thông tin động thời gian thực hiển thị liên tục trong quá trình thám hiểm như thanh sinh mệnh (HP), thanh năng lượng (Mana), thanh thể lực (Stamina), ô phép thuật và bộ đếm tiền vàng. Canvas này sử dụng thành phần \texttt{Canvas Scaler} chế độ \texttt{Scale With Screen Size} với độ phân giải tham chiếu là \texttt{1920x1080} để đảm bảo tỷ lệ co giãn giao diện đồng đều trên các độ phân giải màn hình khác nhau.
    \item \textbf{Menu \& Panels Canvas (Screen Space - Overlay):} Quản lý các giao diện tĩnh xuất hiện khi người chơi tương tác hoặc tạm dừng game (bảng Đồ đạc - Inventory, bảng Trạng thái - Status, bảng Cửa hàng - Shop, bảng Nâng cấp - Facility, Bản đồ toàn cảnh - Full Map). Giao diện này được điều khiển bởi hai lớp trung tâm là \texttt{UIManager} và \texttt{UIStateManager} để kiểm soát vòng đời và độ ưu tiên hiển thị của các bảng chức năng.
    \item \textbf{World Space Canvas:} Sử dụng riêng biệt cho các thành phần UI động cần kết nối trực tiếp với tọa độ thế giới 2D, tiêu biểu là lớp hiển thị chữ sát thương bay \texttt{DamageTextView} xuất hiện ngay trên đầu các thực thể khi nhận sát thương chiến đấu.
\end{itemize}

Để thực hiện cơ chế tạm dừng trò chơi khi mở các bảng chức năng tĩnh (như Main Menu, Status UI, Full Map UI), hệ thống sử dụng lớp \texttt{UIPauseRequester} gửi yêu cầu tới bộ quản lý trung tâm \texttt{GamePauseManager} để thiết lập \texttt{Time.timeScale = 0f}. Việc này dừng hoàn toàn mọi chuyển động vật lý của Rigidbody2D, các script tính toán AI quái vật và hệ thống đạn phép, đảm bảo an toàn cho nhân vật chính trong khi người chơi tương tác với giao diện.

\section{Thiết kế chi tiết các phân hệ}
\label{section:4.6}

Phần này đặc tả chi tiết về cấu trúc thuộc tính, phương thức của các lớp chủ đạo trong hệ thống mã nguồn và cơ cấu lưu trữ dữ liệu bền vững của trò chơi.

\subsection{Thiết kế lớp (Class Design)}
Dưới đây là bảng đặc tả thuộc tính và phương thức của 3 lớp core cốt lõi đóng vai trò xương sống cho trò chơi: \texttt{PlayerController}, \texttt{MvEnemyBase}, và \texttt{GameSaveManager}.

\begin{table}[H]
\centering
\begin{tabular}{|p{2.2cm}|p{4.8cm}|p{8.0cm}|}
\hline
\textbf{Thành phần} & \textbf{Tên Thuộc tính/Phương thức} & \textbf{Vai trò và Ý nghĩa chức năng} \\ \hline
\multirow{8}{*}{\textbf{Thuộc tính}} & \texttt{stateMachine} & Đối tượng máy trạng thái quản lý các state di chuyển của người chơi. \\ \cline{2-3}
 & \texttt{playerInput} & Nhận và lưu trữ đệm các tín hiệu nhập liệu bàn phím/chuột. \\ \cline{2-3}
 & \texttt{playerMovement} & Xử lý các tương tác vật lý (di chuyển, lướt, nhảy, bám tường). \\ \cline{2-3}
 & \texttt{playerStats} & Quản lý các chỉ số runtime sinh mệnh hiện tại của nhân vật. \\ \cline{2-3}
 & \texttt{playerFormManager} & Quản lý việc nạp và hủy các prefab hình thái linh hồn động. \\ \cline{2-3}
 & \texttt{currentForm} & Enum xác định hình thái hiện tại (Dạng người hay biến hình). \\ \cline{2-3}
 & \texttt{currentFormId} & Enum định danh cụ thể hình thái linh hồn đang mở khóa. \\ \hline
\multirow{5}{*}{\textbf{Phương thức}} & \texttt{TakeDamage()} & Triển khai giao diện nhận sát thương từ các nguồn đòn đánh kẻ địch. \\ \cline{2-3}
 & \texttt{Respawn()} & Đưa nhân vật hồi sinh về vị trí Checkpoint đã ghi nhận. \\ \cline{2-3}
 & \texttt{ToggleTransform-}\newline\texttt{Form()} & Kích hoạt/hủy bỏ cơ chế biến đổi linh hồn tại xác quái vật. \\ \cline{2-3}
 & \texttt{AbandonTo-}\newline\texttt{Shrine()} & Đưa nhân vật quay về đền thờ gần nhất khi lựa chọn đầu hàng. \\ \cline{2-3}
 & \texttt{SyncPlayerMovement-}\newline\texttt{Collider()} & Đồng bộ hóa kích thước collider vật lý khớp với hình thái linh hồn mới. \\ \hline
\end{tabular}
\caption{Đặc tả chi tiết thuộc tính và phương thức của lớp PlayerController}
\label{table:player_controller_details}
\end{table}

\begin{table}[H]
\centering
\begin{tabular}{|p{2.2cm}|p{4.8cm}|p{8.0cm}|}
\hline
\textbf{Thành phần} & \textbf{Tên Thuộc tính/Phương thức} & \textbf{Vai trò và Ý nghĩa chức năng} \\ \hline
\multirow{7}{*}{\textbf{Thuộc tính}} & \texttt{maxHealth} & Lượng máu tối đa của quái vật hoặc Boss. \\ \cline{2-3}
 & \texttt{moveSpeed} & Tốc độ di chuyển tuần tra cơ bản. \\ \cline{2-3}
 & \texttt{enemyState-}\newline\texttt{Machine} & Máy trạng thái quản lý các hành vi AI của quái vật. \\ \cline{2-3}
 & \texttt{target} & Tham chiếu Transform của người chơi đang bị truy đuổi. \\ \cline{2-3}
 & \texttt{searchRadius} & Bán kính quét phát hiện người chơi của thực thể. \\ \cline{2-3}
 & \texttt{searchFanAngle} & Góc quét hình quạt phía trước mặt quái vật. \\ \cline{2-3}
 & \texttt{isDead} & Trạng thái xác định quái vật đã bị tiêu diệt hay chưa. \\ \hline
\multirow{6}{*}{\textbf{Phương thức}} & \texttt{TakeDamage()} & Trừ máu hiện tại, thực hiện đẩy lùi và chuyển sang Hit State. \\ \cline{2-3}
 & \texttt{Die()} & Tắt các va chạm collider, kích hoạt VFX tan biến và phát sự kiện chết. \\ \cline{2-3}
 & \texttt{UpdateSearch-}\newline\texttt{Target()} & Thực hiện quét OverlapCircle để xác định mục tiêu tối ưu. \\ \cline{2-3}
 & \texttt{MovePatrol()} & Di chuyển tuần tra qua lại quanh điểm xuất phát ban đầu. \\ \cline{2-3}
 & \texttt{MoveReturnTo-}\newline\texttt{Origin()} & Di chuyển quay lại điểm xuất phát khi mất dấu mục tiêu. \\ \cline{2-3}
 & \texttt{ApplyKnockback()} & Tính toán góc và tác dụng lực vật lý đẩy lùi lên Rigidbody2D. \\ \hline
\end{tabular}
\caption{Đặc tả chi tiết thuộc tính và phương thức của lớp MvEnemyBase}
\label{table:enemy_base_details}
\end{table}

\begin{table}[H]
\centering
\begin{tabular}{|p{2.2cm}|p{4.8cm}|p{8.0cm}|}
\hline
\textbf{Thành phần} & \textbf{Tên Thuộc tính/Phương thức} & \textbf{Vai trò và Ý nghĩa chức năng} \\ \hline
\multirow{6}{*}{\textbf{Thuộc tính}} & \texttt{instance} & Tham chiếu tĩnh duy nhất (Singleton) của trình quản lý. \\ \cline{2-3}
 & \texttt{saveFileName} & Tên mặc định của tệp tin lưu trữ dữ liệu JSON. \\ \cline{2-3}
 & \texttt{activeSlotIndex} & Chỉ mục ô nhớ lưu trữ (Slot Index) đang được tải. \\ \cline{2-3}
 & \texttt{currencyWallet} & Dữ liệu ScriptableObject quản lý ví tiền vàng người chơi. \\ \cline{2-3}
 & \texttt{inventoryState} & Dữ liệu ScriptableObject quản lý túi đồ runtime. \\ \cline{2-3}
 & \texttt{facilityManager} & Tham chiếu quản lý tiến trình nâng cấp vĩnh viễn của làng. \\ \hline
\multirow{6}{*}{\textbf{Phương thức}} & \texttt{SaveActiveSlot()} & Tự động ghi dữ liệu của ô nhớ hiện tại xuống tệp tương ứng. \\ \cline{2-3}
 & \texttt{SaveGameTo-}\newline\texttt{Path()} & Tuần tự hóa đối tượng lưu trữ sang JSON và ghi xuống đĩa cứng. \\ \cline{2-3}
 & \texttt{LoadGameFrom-}\newline\texttt{Slot()} & Đọc tệp JSON tương ứng, giải tuần tự hóa và nạp vào game. \\ \cline{2-3}
 & \texttt{CreateNew-}\newline\texttt{GameSlot()} & Khởi tạo ô nhớ mới với các dữ liệu mặc định ban đầu. \\ \cline{2-3}
 & \texttt{CaptureSave-}\newline\texttt{Data()} & Tổng hợp toàn bộ dữ liệu từ các phân hệ vào đối tượng GameSaveData. \\ \cline{2-3}
 & \texttt{ApplySave-}\newline\texttt{Data()} & Phân bổ dữ liệu từ đối tượng GameSaveData nạp ngược vào hệ thống. \\ \hline
\end{tabular}
\caption{Đặc tả chi tiết thuộc tính và phương thức của lớp GameSaveManager}
\label{table:save_manager_details}
\end{table}

Để minh họa luồng truyền thông điệp giữa các đối tượng trong các tình huống tương tác thực tế, hệ thống đặc tả 2 use case cốt lõi thông qua biểu đồ trình tự:
\begin{enumerate}
    \item \textbf{Use Case Chuyển đổi linh hồn (Form Switch):} Sơ đồ trình tự tương tác giữa \texttt{PlayerInput}, \texttt{PlayerController}, \texttt{PlayerFormManager}, và \texttt{PlayerStats} khi thực hiện thay đổi hình dạng nhân vật được biểu diễn chi tiết tại Hình~\ref{fig:form_switch_sequence} (Mục~\ref{section:4.2}).
    \item \textbf{Use Case Lưu trữ tiến trình game (Save Game):} Khi người chơi tương tác cầu nguyện tại đền thờ hồi sinh (\texttt{RespawnShrine}), một yêu cầu lưu game được gửi tới \texttt{GameSaveManager}. Bộ quản lý sẽ gọi hàm \texttt{CaptureSaveData()} để thu thập dữ liệu từ các ScriptableObjects trạng thái (\texttt{CurrencyWalletSO}, \texttt{InventoryStateSO}, v.v.) và các service tiến trình khác, biên dịch thành cấu trúc trung gian \texttt{GameSaveData}. Đối tượng này được chuyển đổi thành chuỗi JSON qua bộ tiện ích của Unity và ghi xuống đĩa cứng dưới dạng tệp tin văn bản cục bộ.
\end{enumerate}

\subsection{Thiết kế cơ sở dữ liệu (Database Design)}
Vì trò chơi được thiết kế như một phần mềm độc lập chạy ngoại tuyến hoàn toàn trên PC (PC Standalone), hệ thống không sử dụng các hệ quản trị cơ sở dữ liệu quan hệ máy chủ (như SQL Server hay MySQL) để tránh đòi hỏi kết nối mạng và các chi phí cài đặt phức tạp cho người chơi.

Thay vào đó, trò chơi áp dụng thiết kế cơ sở dữ liệu hướng tài liệu cục bộ (document-based local file database) bằng định dạng JSON bền vững. Cơ sở dữ liệu được tổ chức dưới dạng một lớp tổng hợp gốc là \texttt{GameSaveData} (đóng vai trò như một Document) chứa các tập hợp dữ liệu con đại diện cho trạng thái của từng phân hệ (đóng vai trò như các Sub-documents). Cấu trúc chi tiết của các trường dữ liệu trong tài liệu lưu trữ JSON này được đặc tả hoàn chỉnh trong Bảng~\ref{table:json_save_data} (Mục~\ref{section:4.4.1}). Việc tổ chức cơ sở dữ liệu dưới dạng JSON giúp đảm bảo tốc độ đọc ghi cực nhanh, dung lượng lưu trữ nhẹ (dưới 100KB mỗi slot lưu trữ) và dễ dàng sao lưu, phục hồi thông qua các thư viện hệ thống tệp chuẩn của hệ điều hành.

\end{document}
